
\documentclass[12pt]{article}
\usepackage{amsmath} % AMS Math Package
\usepackage{amsthm} % Theorem Formatting
\usepackage{amssymb}    % Math symbols such as \mathbb
\usepackage{graphicx} % Allows for eps images
\usepackage[dvips,letterpaper,margin=1in,bottom=0.7in]{geometry}
\usepackage{tensor}
\newtheorem{p}{Problem}[section]
\usepackage{amsmath}

\begin{document}

{\noindent\Huge\bf  \\[0.5\baselineskip] {\fontfamily{cmr}\selectfont  Problem Set I}         }\\[2\baselineskip] % Title
{ {\bf \fontfamily{cmr}\selectfont Political Science 43401}\\ {\textit{\fontfamily{cmr}\selectfont     October 1, 2020}}}\hfill   {\large \textsc{Robert Gulotty}} % Author name
\\[1.4\baselineskip] 





\begin{enumerate}

\item If sets A and B are disjoint, what do I know about $x$ if I know that $x\in B$?  Give a concrete example.


\item[3.1.1]  For each of the following pairs of sets A and B, state which, if either, is a subset of the other:

(a) A is the set of people living in Scotland; B is the set of people living in the UK.

(b) A is the set of people living in the USA; B is the set of people living in California.

(c) A is the set of natural numbers divisible by 6; B is the set of natural numbers divisible by 2.

(d) A is the set of natural numbers divisible by 5; B is the set of natural numbers divisible by 7.


\item[3.1.3] Let $A, B, C$ be subsets of a set $X$.  Which of the expressions $(A\cap B)\cup (A\cap B^c)$, $(A\cap B)\cup (A\cap B^c)^c$, $(A\cap B)\cup (A^c\cup C)^c$ can be simplified, and how?


\item[3.4.2] Let $g$ be the mapping from $\mathbb{R}^2$ to $\mathbb{R}^2$ defined by $$g(x, y)=(y, -x),$$ and let the mapping $h$ be as in Exercise 3.4.1.  Find the image of $(x, y)$ under each of the composite mappings $g\circ h$ and $h\circ g$.


\item[31.1.1] As in the text, let $\mathbb{Q}$ be the set of all rational numbers. If $x\in \mathbb{Q}$, which of the following statements (P, Q, R, S) imply which of the others?

\begin{itemize}

\item[P:] $x\geq 0$
\item[Q:] $\exists w\in \mathbb{Q} $ such that $x=w^2$
\item[R:] $nx>-1$ $\forall n\in \mathbb{N}$
\item[S:] $\exists n\in \mathbb{N}$ such that $nx>-1$.

\end{itemize}
\item[R Bonus] Using R, plot the following functions:
$$y(x)=ln(x)+4 $$
$$y(x)=\frac{1}{2}e^{x^2/2}$$
$$z(x,\ y)=3x^{.5}y^{.5}$$

\end{enumerate}
\end{document}